\chapter{From Scripts to a Station Pipeline}
\label{ch:course-pipeline}

The preceding chapters introduced DAGonStar's abstractions one at a time. This
chapter assembles them into the reference workflow of the course: the Gulf of
Naples station pipeline of Lecture~12. It also closes Part~I by showing the two
operations that make a workflow an object one can reason about, rather than a
program one can only run: serialisation (Lecture~11) and comparison with the
script it replaces (Lecture~02).

Every listing in this chapter is included verbatim from the course
repository. The command that runs it is given next to it, and its recorded
output follows. The script \texttt{tools/run\_examples.py} re-runs every
example and compares its output with the recorded one, so what the book shows
is what the code does.

\section{The Course Dataset}

All examples in Parts~I and~II share one small dataset: twelve hourly
observations from three coastal stations of the Gulf of Naples --- Naples
(\texttt{NAP}), Capri (\texttt{CAP}) and Ischia (\texttt{ISC}) --- on
2~November~2026. Each record carries a station code, a UTC time, an air
temperature in degrees Celsius and a wind speed in metres per second. The
synthetic dataset is stored in \texttt{examples/data/gulf\_of\_naples.csv}, and
the helpers in \texttt{examples/common/course.py} copy it into each workflow's
\texttt{fetch} task. Every example therefore starts from the same bytes,
without a download, and runs in a fresh temporary scratch directory that it
removes when it ends.

\begin{practicebox}
A teaching dataset should be small enough to check by hand and rich enough
to exercise every mechanism. Four records per station allow a mean to be
verified mentally (for Naples, $62/4 = 15.5$\,\textdegree C), and three stations
are the smallest number for which fan-out and fan-in are visibly different
from a chain.
\end{practicebox}

\section{A Script and a DAG}

Lecture~02 computes the three station means twice: once as a plain Python
script and once as a DAGonStar workflow with one task per station and a
report that waits for all three.

\lstinputlisting[caption={Lecture 02: the same analysis as a script and as a DAG (\texttt{examples/02-script-vs-dag/as\_dag.py}).},label={lst:as-dag}]{../../examples/02-script-vs-dag/as_dag.py}

\begin{lstlisting}[style=bash]
python examples/02-script-vs-dag/as_dag.py
\end{lstlisting}

\lstinputlisting[style=bash,caption={Recorded output of Lecture 02.}]{../../examples/02-script-vs-dag/output/as_dag.txt}

The two versions agree on the numbers. They differ in what each makes explicit. The
script's order is an accident of how its lines were written. The DAG's order
is derived from the \texttt{workflow://} references, so the three means are
independent and may run concurrently, and the report's dependencies are
listed by the engine (\texttt{report waits for: [...]}). The DAG also leaves
behind one folder per task with its launcher and log. The script leaves only
its result.

\section{Inspecting a Workflow as Data}

Lecture~11 treats a workflow as a document. \texttt{as\_json()} produces a
dictionary with the keys \texttt{name}, \texttt{id}, \texttt{host} and
\texttt{tasks}, and a new \texttt{Workflow} can be rebuilt from it with
\texttt{jsonload}. Timing and lifecycle events are read from the same
object.

\lstinputlisting[caption={Lecture 11: serialise, reload, time and observe (\texttt{examples/11-inspect/inspect\_workflow.py}).},label={lst:inspect}]{../../examples/11-inspect/inspect_workflow.py}

\lstinputlisting[style=bash,caption={Recorded output of Lecture 11.}]{../../examples/11-inspect/output/inspect_workflow.txt}

\begin{definitionbox}{Explicit edges in the JSON form}
A reference inside a command is rediscovered when the reloaded workflow
parses its commands again. An edge added with \texttt{add\_dependency\_to()}
has no textual trace in any command. It survives the round trip only because
\texttt{as\_json()} records it in \texttt{explicit\_prevs} and
\texttt{make\_dependencies()} restores it. This is why the reloaded
\texttt{plot} still waits for \texttt{count}.
\end{definitionbox}

\section{The Station Pipeline}

The pipeline of Lecture~12 brings together every mechanism of Part~I. A
\texttt{fetch} task produces the observations. Three station tasks compute
the mean temperature and the maximum wind of each station, and a
\texttt{report} task builds a table and names the windiest station.

\lstinputlisting[caption={Lecture 12: the station pipeline (\texttt{examples/12-station-pipeline/station\_pipeline.py}).},label={lst:station-pipeline}]{../../examples/12-station-pipeline/station_pipeline.py}

\begin{lstlisting}[style=bash]
python examples/12-station-pipeline/station_pipeline.py
\end{lstlisting}

\lstinputlisting[style=bash,caption={Recorded output of Lecture 12.}]{../../examples/12-station-pipeline/output/station_pipeline.txt}

\begin{figure}[h]
\centering
\begin{tikzpicture}[
  node distance=10mm and 16mm,
  task/.style={rectangle, rounded corners, draw=dagonblue, very thick,
    fill=blue!4, minimum width=24mm, minimum height=9mm, font=\sffamily\small},
  arr/.style={-{Stealth[length=2.4mm]}, thick, dagonblue}
]
  \node[task] (fetch) {fetch};
  \node[task, right=of fetch] (cap) {summary\_CAP};
  \node[task, above=of cap] (nap) {summary\_NAP};
  \node[task, below=of cap] (isc) {summary\_ISC};
  \node[task, right=of cap] (report) {report};
  \foreach \s in {nap,cap,isc} { \draw[arr] (fetch) -- (\s); \draw[arr] (\s) -- (report); }
\end{tikzpicture}
\caption{The station pipeline: one producer, three independent consumers, one
fan-in. Every edge is inferred from a \texttt{workflow://} reference.}
\label{fig:station-pipeline}
\end{figure}

\subsection{Reading the Result}

The table can be checked against the dataset by hand. Capri is the warmest
station (16.75\,\textdegree C) and the windiest (8.0\,m/s). Naples is the
coolest (15.50\,\textdegree C). The last line reports that five tasks ended,
which matches the five nodes of \Cref{fig:station-pipeline}. A workflow that
silently skipped a station would report four.

\section{What Part~I Has Established}

\begin{itemize}[leftmargin=1.4em]
  \item A workflow is a graph whose edges come from data references
    (\Cref{ch:schema}) or explicit calls (\Cref{ch:lessons-01-02}).
  \item Every task leaves a folder with a launcher, a log and its outputs
    (\Cref{ch:lessons-03-04}).
  \item The graph is validated before it runs, and cycles are reported with
    the task involved (\Cref{ch:lessons-05-07}).
  \item A workflow can be serialised, reloaded, timed and observed, so it can
    be studied as data.
\end{itemize}

\begin{exercise}
Add a fourth station, Pozzuoli (\texttt{POZ}), to the dataset helper and to
the pipeline. Predict the new table before running the example, run it, and
explain any difference between prediction and result.
\end{exercise}
